\chapter{प्रस्तावना}

इस पुस्तक की महत्वाकांक्षा सरल है: पृथ्वी और अंतरिक्ष विज्ञान की एक ही सुसंगत पुस्तक, जो
विद्यालय के पहले वर्ष से स्नातक की डिग्री के अंत तक सब कुछ समेटे, दुनिया में कहीं के भी
पाठकों के लिए लिखी गई, पहले अंग्रेज़ी में, और यहाँ हिंदी में अनूदित।

शैली संक्षिप्त और कसी हुई है: परिभाषाओं, उदाहरणों, प्रतिज्ञप्तियों, विधियों और आरेखों से
बने पाठ, जहाँ भी दिए गए स्तर पर संभव हो वहाँ सावधान तर्क के साथ, और उनके बाद ध्यान से चुने
गए अभ्यास। जो परिणाम पूरी व्युत्पत्ति के बिना दिए गए हैं, उन्हें स्पष्ट रूप से स्वीकृत के रूप
में चिह्नित किया गया है। सभी अभ्यासों के पूरे हल पुस्तक के अंत में इकट्ठे दिए गए हैं। असली
पृथ्वी, उसके महासागरों, उसके वायुमंडल या आकाश के बारे में दी गई हर संख्या किसी संदर्भ स्रोत
से आती है, और असली मापों से बनाया गया हर वक्र प्रकाशित आँकड़ों से निकाला गया है।

पाठ्यक्रम पाँच पुस्तकों में है। पहली प्राथमिक और माध्यमिक विद्यालय के वर्षों को समेटती है,
दूसरी उच्च माध्यमिक के तीन वर्षों को, और आख़िरी तीन विश्वविद्यालय के तीन वर्षों को। विद्यालय
की दोनों पुस्तकें कोई भी धारणा कभी दो बार नहीं पढ़ातीं: हर विचार एक ही बार पढ़ाया जाता है,
उस वर्ष में जहाँ वह पहली बार ठीक बैठता है, और बाद का कोई अध्याय जिसे उसकी ज़रूरत हो, आगे
बढ़ने से पहले एक छोटे-से अनुस्मारक, \emph{पहले से ज्ञात}, से शुरू होता है। जहाँ कोई धारणा
भौतिकी, जीव विज्ञान या रसायन विज्ञान की है, वहाँ उसी संग्रह की साथी पुस्तकें उसे पढ़ाती हैं,
और यह पुस्तक पृथ्वी वैज्ञानिक का दृष्टिकोण अपनाने से पहले उसे याद दिलाती है। विद्यालय की
पुस्तकों में हर वर्ष एक भाग है, जो अध्यायों में बँटा है; विश्वविद्यालय की हर पुस्तक एक वर्ष है,
जो अध्यायों में बँटी है। वर्ष जितना शुरुआती, पाठक उतना ही छोटा जिसके लिए वह लिखा गया है:
ज़्यादा चित्र, ज़्यादा विस्तार से हल किए गए चरण, और ज़्यादा सरल अभ्यास।

\section*{यह पुस्तक कैसे पढ़ें}

कथन हर अध्याय के भीतर क्रमांकित हैं और रंगों से पहचाने जाते हैं: परिभाषाएँ नीले में, प्रमेय
लाल में, प्रतिज्ञप्तियाँ और उनके संबंधी नारंगी में, और विधियाँ --- सामान्य कामों के लिए छोटे
नुस्ख़े --- हरे में। अभ्यासों को सीधे अनुप्रयोगों के लिए ($\star$) से लेकर कठिन समस्याओं के
लिए ($\star\star\star$) तक श्रेणीबद्ध किया गया है। फ़्रेम वाले बक्सों में वह है जो तर्क का
हिस्सा नहीं है: किसी पिछले अध्याय में पढ़ाई गई बात का अनुस्मारक, भूवैज्ञानिक और दूसरे पृथ्वी
वैज्ञानिक क्षेत्र-कार्य में या प्रयोगशाला में कैसे काम करते हैं, इतिहास का एक पन्ना, वर्ष-दर-वर्ष
बदलने वाले आँकड़े (उस वर्ष के साथ जिसके लिए वे दिए गए हैं), और सुरक्षा की सलाह।

\section*{सुरक्षा}

पृथ्वी को उसे देखकर सीखा जाता है: आकाश, मौसम, कोई भूदृश्य, किसी इमारत के पत्थर, किसी बगीचे
की मिट्टी। यह पुस्तक अपने पाठक से बस इतना ही माँगती है। सूर्य की ओर कभी मत देखो, न किसी
उपकरण के साथ न उसके बिना: वह अंधा कर देता है। खदानें, कगारें, गुफाएँ, उफनती नदियाँ और
ज्वालामुखी प्रशिक्षण और उपकरणों वाले पेशेवरों की जगहें हैं, और यह पुस्तक उनके काम का वर्णन
करती है, किसी को वहाँ भेजे बिना।

\section*{योगदान}

इस पुस्तक का स्रोत एक सार्वजनिक git भंडार में है:
\begin{center}
\url{https://github.com/bvirrion/one-earth-and-space-sciences-book}
\end{center}
योगदान --- नए अध्याय, सुधार, बेहतर व्याख्याएँ, अतिरिक्त अभ्यास --- का स्वागत है; भंडार के
मूल में \texttt{CONTRIBUTING.md} देखो।
